\documentclass[10pt]{article}

\usepackage[final]{neurips_2019}
\usepackage[top=0.7in, bottom=0.7in, left=0.75in, right=0.75in]{geometry}

\usepackage[utf8]{inputenc}
\usepackage[T1]{fontenc}
\usepackage{hyperref}
\usepackage{url}
\usepackage{booktabs}
\usepackage{amsfonts}
\usepackage{nicefrac}
\usepackage{microtype}
\usepackage{graphicx}
\usepackage{xcolor}
\usepackage{lipsum}
\usepackage{enumitem}
\usepackage{amsmath}
\usepackage{colortbl}
\usepackage{float}
\usepackage{fancyhdr}



\newcommand{\note}[1]{\textcolor{blue}{{#1}}}

% Remove headers and footers
\pagestyle{empty}
\fancyhf{}
\renewcommand{\headrulewidth}{0pt}
\renewcommand{\footrulewidth}{0pt}

\title{
  Layer-Adaptive MX Format Quantization for RLHF Energy Optimization on FPGAs \\
  \vspace{0.5em}
  \small{\normalfont Stanford CS217 Project -- Progress Report}  % Progress Report
}

\author{
  Hiva Zaad (Mohammadzadeh) \\
  Department of Computer Science \\
  Stanford University \\
  \texttt{hiva@stanford.edu} \\
  \And
  Grant Griffith \\
  Graduate School of Business \\
  Stanford University \\
  \texttt{grgriff@stanford.edu} \\
  \And
  Daniel Adkins \\
 Graduate School of Business \\
  Stanford University \\
  \texttt{dadkins@stanford.edu}
}

\begin{document}

\maketitle
\vspace{-1em}

\begin{abstract}
\small
This project investigates whether layer-adaptive Microscaling (MX) format quantization can reduce the energy consumption of matrix multiply (GEMM) operations during RLHF training on an FPGA, compared to an INT8 baseline. The model is Qwen2.5-0.5B trained with PPO on 1,200 Anthropic HH-RLHF preference pairs (1000 train / 200 eval). All GEMM operations are offloaded to an AWS F2 FPGA (Xilinx VU9P) via AXI interface; the host CPU handles the full training loop.
\end{abstract}
\vspace{-0.5em}

\section{Milestone Status}
\vspace{-0.3em}

\begin{table}[H]
\raggedright
\footnotesize
\setlength{\tabcolsep}{4pt}
\renewcommand{\arraystretch}{1.1}
\begin{tabular}{@{}p{3cm}p{1.2cm}p{1.6cm}p{7.5cm}@{}}
\toprule
\textbf{Milestone} & \textbf{Owner} & \textbf{Status} & \textbf{Key Deliverables} \\
\midrule
\textbf{M1: Environment Setup (Week 1)} & All & \cellcolor{green!30}\textbf{COMPLETE} &
Repo structured. Lab 1 INT8 PE core forked. AWS F2 FPGA access confirmed. Qwen2.5-0.5B validated. Energy protocol: Power × Runtime via XPE.

\textbf{Github:} \scriptsize\url{https://github.com/HivaMohammadzadeh1/CS217-Final-Project/tree/main}

\textbf{Dataset:} \scriptsize\url{https://huggingface.co/datasets/hivamoh/cs217-rlhf-dataset} \\
\midrule
\textbf{M2.1: INT8 FPGA Baseline Energy (Week 2)} & All & \cellcolor{green!30}\textbf{COMPLETE} &
RLHF loop with FPGA matmul offload. 50 PPO steps on 1000 pairs (Qwen2.5-0.5B + reward model). Per-phase energy logged. XPE power estimates extracted, baseline CSV generated. \\
\midrule
\textbf{M2.2: PyTorch MX Sensitivity Profiling (Week 2)} & Grant & \cellcolor{orange!50}\textbf{IN PROGRESS} &
mx-pytorch installed. Per-layer sweep: MXFP4/MXFP8 on transformer layers with group sizes 8, 16. Perplexity deltas logged. Sensitivity matrix and policies (A--D) to be finalized. \\
\midrule
\textbf{M3: SystemC MX Datapath (Week 3)} & Hiva & \cellcolor{gray!30}\textbf{UPCOMING} &
Dual-precision MXFP4/MXFP8 PE in SystemC. Testbenches vs PyTorch reference. \\
\midrule
\textbf{M4: FPGA Synthesis + Integration (Week 4)} & Hiva + Daniel & \cellcolor{gray!30}\textbf{UPCOMING} &
HLS synthesis. Bitstream generation and AWS F2 deployment. End-to-end validation. \\
\midrule
\textbf{M5: Energy Experiments (Week 5)} & All & \cellcolor{gray!30}\textbf{UPCOMING} &
Run MX policies (A--D) on FPGA. Compare energy vs INT8. Generate Pareto curves. \\
\midrule
\textbf{M6: Analysis + Report (Week 6)} & All & \cellcolor{gray!30}\textbf{UPCOMING} &
Final report with Pareto analysis, policy recommendation, and MX viability conclusions. \\
\bottomrule
\end{tabular}
\end{table}
\vspace{-0.5em}

\section{Milestone 2.1 Detail: INT8 FPGA Baseline Energy}
\vspace{-0.3em}

\small
The baseline energy measurement establishes the reference point for all subsequent MX format comparisons. Every claimed energy saving in the final report is expressed as a delta from this number.

\subsection{What Was Built}
\vspace{-0.2em}

\begin{itemize}[leftmargin=*, noitemsep, topsep=0pt]
\item \textbf{Tiling + offload:} Breaks matmuls into 16×16 tiles, sends to FPGA via AXI.
\item \textbf{RLHF script:} Loads Qwen2.5-0.5B + reward model, hooks matmul offload into PPO loop.
\item \textbf{Baseline run:} 50 PPO steps on 1000 pairs. Runtime logged by phase.
\item \textbf{Energy:} XPE power × Runtime for each phase, written to CSV.
\end{itemize}
\vspace{-0.2em}

\subsection{Output Artifact}
\vspace{-0.2em}
Baseline CSV: layer\_type, rlhf\_phase, avg\_power\_W, runtime\_s, energy\_J.
\vspace{-0.2em}

\subsection{Dataset}
\vspace{-0.2em}
HuggingFace: hivamoh/cs217-rlhf-dataset. 1,200 pairs (1000/200 split), 512 tokens max.
\vspace{-0.2em}

\subsection{Next Steps (Week 2)}
\vspace{-0.2em}

\begin{itemize}[leftmargin=*, noitemsep, topsep=0pt]
\item Finalize per-layer sweep MXFP4/MXFP8/INT8. Define 2\% perplexity tolerance.
\item Produce sensitivity matrix: layers × format × group size vs. INT8 baseline.
\item Define MX policies (Conservative, Balanced, Aggressive, Phase-Adaptive) for Week 5.
\item Hand off sensitivity results for Week 3 SystemC datapath design.
\end{itemize}
\vspace{0.3em}

\noindent\scriptsize{GitHub: github.com/HivaMohammadzadeh1/CS217-Final-Project | Dataset: huggingface.co/datasets/hivamoh/cs217-rlhf-dataset}

\end{document}
